% This is the Introduction to the Hodgkin-Huxley neuron model V2.0 book
% Last edited 2026.08.09

\ifx\magyar\undefined
\chapter{Introduction\label{ch:Introduction}}
\else
\chapter{Bevezetés\label{ch:Introduction}}
\fi
\iflatexml
\else
\minitoc
\fi


\quotationbox
{"The Human Brain Project should lay the technical foundation for a \textbf{new model} of ICT-based brain research, driving \textit{integration between data and knowledge from different disciplines}, and catalysing a community effort to achieve a \textbf{new understanding of the brain}%, new treatments for brain disease 
\dots and \textit{new brain-like computing technologies}."\\
{the \href{https://www.humanbrainproject.eu/en/}{Human Brain Project},  summarised its goal @2012}}

Az egyes neuronok \textit{dinamikus} működése, kapcsolataik,
magasabb szintű szer\-veződéseik és azok kapcsolatai;
az agy információ feldolgozási képessége; végül 
a tudat és viselkedés megértése még ma is a tudomány 
nagy rejtélyei közé tartoznak:
\hyperlink{LivingMatter}
{mely ponton kel életre az élettelen anyag}, \textit{\hyperlink{CognitiveMatter}{mitől válik intelligenssé}}; továbbá \hyperlink{PhysicalPictureNeuron}{hogyan tudja ezeket a kérdéseket a tudomány megválaszolni}. 
Valóban \hyperlink{NewUnderstanding}
{másfajta megértésre van szükség}; ezt nyújtjuk a könyvben.
%
Az újonnan indított projektek iránti várakozás óriási, 
de úgy tűnik, mintha a szilárd tudományos megalapozás reménye
elhalványuult volna.
Például, az újonnan (2025 végén) "\href{https://brainminds.jp/en}{Brain/MINDS 2.0}" Japánban indított program  már csak 
matematikai modelleket használ, \textit{meg sem próbálja megérteni az alapul szolgáló fizikai folyamatokat}.
A neuron model\-lekben a 'dinamikus' arra utal, hogy a neuron
állapotváltozói (elsősorban a membrán potenciál és az ion\-csatornák tulajdonságai) hogyan viselkednek az idő függvényében belső folyamatok és külső körülmények hatására.
A mi model\-lünktől eltérően, ezek a folyamatok csupán egy 
idő paramétert használnak, egy empírikus vagy egy matematikai függvényben, lásd pl.~\cite{NeuralDynamicsGertsner:2014}.
Mi valódi, fizikai alapokon nyugvó dinamikát használunk a neuronok reakciójának leírására, olyan értelemben, ahogyan Newton vezette be mozgástörvényeit.

\index{laws of motion}

\section{Tudományterületi rétegződés\label{sec:Intro-Disciplinary}}



A könyv több rétegű. A mindenki (azaz elsősorban a megfelelő
fiziológiai háttérrel rendelkező olvasó)
által megérthetőnek szánt részek mellett a vonatkozó más jellegű háttérrel rendelkezők számára találhatók 
a szükséges
\textbf{\hyperlink{intro:MyMathematicsPanel}{matematikai}} és
\textbf{\hyperlink{intro:MyPhysicsPanel}{fizikai}}
ismereteket tartalmazó panelek is, amelyeknek megértéséhez legalább a megfelelő szintű BSc ismeretek, de néha ezt túllépő tudás, esetenként 
a gondolatok szokatlanságával elmélyülésre késztető 
ismeretek és magyarázatok találhatók.
A jelenségkör bonyolultsága sok közelítést tesz szükségessé,
de a megfelelő közelítések és absztrakciók felhasználásával 
numerikus eredményeket (néha csupán észszerű becsléseket)
tartalmaznak a \textbf{\hyperlink{num:MyNumericsPanel}{numerikus}} panelek.
Törekszünk arra, hogy az anyag ezek nélkül is érthető legyen,
ezekkel együtt pedig a megfelelő szakképzettségű olvasó
továbbgondolásra ösztönző tudást meríthessen.



\physicspanel{intro:MyPhysicsPanel}{Fizikai panel, példa}
{
A tudományok specializálódása során a fizika az élettelen természet 
leírására szakosodott és (az észszerű leírhatóság érdekében)
korlátozott érvényességű szakterületeket hozott létre.
Az élő természet tanulmányozása során már természetes módon
használták eredményeit, módszereit és eszközeit.
Ennek során abból a helyes nézőpontból indultak ki, hogy
a természetben található objektumok és kölcsönhatásaik leírására
ugyanazoknak a törvényeknek kell vonatkozni. Nem vették azonban figyelembe, hogy a természet bonyolultsága miatt,
a leírhatóság érdekében, közelítéseket és elhanyagolásokat kell tennünk. A diszciplináris törvények egy-egy jelenségköre vonatkozóan
absztrakt fogalmakat alkotnak és a törvények ezeket az absztrakt
fogalmakat írják le pontosan, nem pedig azokat a jelenségeket,
amelyek a természetben előfordulnak. Az esetek jelentős részében
ezek a törvények teljes egészükben használhatók, 
az élő és élettelen környezet körülményeinek jelentős eltérései miatt
sokszor csak jelentős módosítással vagy akár teljes újragondolással alkalmazhatók; lásd Erwin Schrödinger lényeglátó kérdéseinek értelmezését
\index{Schrödinger, Erwin}
a \textbf{Physics Panel~\ref{phys:SchrödingerNonOrdinaryLaws}} 
}



\mathematicspanel{intro:MyMathematicsPanel}{Matematikai panel, példa}
{
A matematikai panelek egyenleteket és azok értelmezését tartalmazzák. A fő szöveg ezek nélkül is érthető állításokat tartalmaz, az egyenletek -- a kellő matematikai ismeretekkel
is rendelkező olvasók számára -- ezeket támasztják alá.
}


\numericspanel{intro:MyNumericsPanel}{Numerikus panel, példa}
{
A tapasztalatok összetettsége sok  közelítést tesz szükségessé.
Megfelelő közelítések és absztrakciók alkalmazásával
numerikus eredményeket (bár néha csupán észszerű becsléseket)
tudunk megadni a tárgyalt jelenségekre.
}


\quotationbox
{
"Biological laws are much more difficult to find than physical ones"~\cite{VollmerLimitsOfBiology:1995} @1995
}




\section[Abstract discussion]{Abstract discussion\label{sec:Intro-Abstract}}

Nature uses an infinite variety of implementing neurons. 
Our goal is not to discuss all
infinitely complex molecular, biochemical and physiological details
of neuronal operation
In the \gls{CNS} 
they can cooperate
with each other, 'despite the extraordinary diversity and complexity of neuronal morphology and synaptic connectivity,
\textit{the nervous systems adopt \textbf{a number of
 basic principles}}
for all neurons and synapses',
\cite{JohnstonWuNeurophysiology:1995}
independently from the infinitely complex molecular, biochemical and physiological details.
\index{Johnston Daniel}
\index{Wu Samuel Miao-sin }
We \textit{will base our holistic discussion on those general basic principles 
and create an 
\hyperlink{NeuronPhysical}
{'abstract physical neuron'}; a hypothetical element, which functions
essentially like a neuron},
skipping the 'implementation details' nature uses.
We essentially follow von Neumann's method in describing the foundations of
computing science~\cite{EDVACreport1945}: "we will base our considerations on a hypothetical element, which functions essentially like a neuron".
We need \hyperlink{Abstractions}
{different abstractions and approximations} 
(furthermore, \hyperlink{NonOrdinaryLaws}{'non-ordinary' laws} of science) for describing biological processes.
\index{ordinary laws}
However, an abstraction is usable
in practice only when paired with a generalization: the more abstract the assumption, the more general and widely applicable the concept or conclusion.
\textit{By abstractions, we can reduce the unbelievably detailed world into manageable pieces, and we can learn anything general.}
We must show which approximations are oversimplifications and which phenomena are misunderstood; measured, or interpreted in a wrong approach.
Abstraction is, in fact, everywhere, including inside each of us. It is a core element of science and cognition. 

